\paragraph{冻结微正则分解与原子 residue 的单维扰动}
在冻结 escort 态向最大纤维层的指数集中、以及单原子 Euler 因子在 periodic/core residue 侧的精确剥离都已建立之后，还可把这两条链进一步压缩为一组完全闭式的终端结论：冻结相中的 escort 态相对于基态微正则分布具有精确二分解，而原子 residue 修正在有限 Fourier 侧则形成单维等幅相位向量。

\begin{conclusion}[冻结相 escort 态相对于基态微正则分布的精确二分解与 \texorpdfstring{$f$}{f}-散度闭式]\label{con:xi-time-part57ba-freezing-escort-groundstate-fdiv-closed}
沿用定理 \ref{thm:xi-fixed-freezing-escort-renyi-spectrum-collapse} 与推论 \ref{cor:xi-fixed-freezing-escort-groundstate-tv-identity} 的记号，固定 \(a>a_{\mathrm c}\)。记
\[
\sigma_m:=\mathrm{Unif}(A_m^{\max})
=\frac{1}{K_m}\ind_{A_m^{\max}}(\cdot),
\qquad
p_{m,a}:=\pi_{m,a}(A_m^{\max}).
\]
则存在概率测度 \(\eta_{m,a}\)；当 \(p_{m,a}<1\) 时其支撑于 \(X_m\setminus A_m^{\max}\)，并且
\[
\pi_{m,a}=p_{m,a}\sigma_m+(1-p_{m,a})\eta_{m,a},
\qquad
\pi_{m,a}(\,\cdot\mid A_m^{\max})=\sigma_m.
\]
进一步，对任意满足 \(f(1)=0\) 且 \(f(0)<\infty\) 的 Csisz\'ar--Morimoto 散度
\[
D_f(\sigma_m\|\pi_{m,a})
:=
\sum_{x\in X_m}\pi_{m,a}(x)\,
f\!\left(\frac{\sigma_m(x)}{\pi_{m,a}(x)}\right),
\]
都有精确闭式
\[
\boxed{\;
D_f(\sigma_m\|\pi_{m,a})
=
p_{m,a}\,f\!\left(\frac{1}{p_{m,a}}\right)
+(1-p_{m,a})f(0).\;}
\]
特别地，
\[
\boxed{\;
\|\pi_{m,a}-\sigma_m\|_{\mathrm{TV}}=1-p_{m,a},\;}
\]
\[
\boxed{\;
D_{\mathrm{KL}}(\sigma_m\|\pi_{m,a})=-\log p_{m,a},\;}
\]
\[
\boxed{\;
\chi^2(\sigma_m\|\pi_{m,a})=\frac{1}{p_{m,a}}-1.\;}
\]
并且对任意固定的 \(f\)，上述散度满足
\[
D_f(\sigma_m\|\pi_{m,a})
=
O\!\bigl(\exp(-m\Delta(a)+o(m))\bigr).
\]
\end{conclusion}
\begin{proof}
由于 \(d_m(x)=M_m\) 在 \(A_m^{\max}\) 上恒定，对任意 \(x\in A_m^{\max}\) 都有
\[
\pi_{m,a}(x)=\frac{M_m^a}{S_a(m)}=\frac{p_{m,a}}{K_m}.
\]
因此 \(\pi_{m,a}\) 在 \(A_m^{\max}\) 上的条件分布恰为 \(\sigma_m\)。若 \(p_{m,a}<1\)，取 \(\eta_{m,a}\) 为 \(\pi_{m,a}\) 在补集 \(X_m\setminus A_m^{\max}\) 上的条件分布；若 \(p_{m,a}=1\)，由于第二项系数为零，任取 \(X_m\) 上概率测度 \(\eta_{m,a}\) 即可。于是得到所述精确分解。

再计算 \(f\)-散度。对 \(x\in A_m^{\max}\)，有
\[
\frac{\sigma_m(x)}{\pi_{m,a}(x)}
=
\frac{1/K_m}{p_{m,a}/K_m}
=
\frac{1}{p_{m,a}},
\]
而对 \(x\notin A_m^{\max}\)，有 \(\sigma_m(x)=0\)。故
\[
\begin{aligned}
D_f(\sigma_m\|\pi_{m,a})
&=
\sum_{x\in A_m^{\max}}\pi_{m,a}(x)\,
f\!\left(\frac{1}{p_{m,a}}\right)
+
\sum_{x\notin A_m^{\max}}\pi_{m,a}(x)\,f(0)\\
&=
p_{m,a}\,f\!\left(\frac{1}{p_{m,a}}\right)
+(1-p_{m,a})f(0).
\end{aligned}
\]
这给出一般闭式。

对 \(f(t)=t\log t\)，得到
\[
D_{\mathrm{KL}}(\sigma_m\|\pi_{m,a})
=
p_{m,a}\cdot \frac{1}{p_{m,a}}\log\frac{1}{p_{m,a}}
=
-\log p_{m,a}.
\]
对 \(f(t)=(t-1)^2\)，得到
\[
\chi^2(\sigma_m\|\pi_{m,a})
=
p_{m,a}\left(\frac{1}{p_{m,a}}-1\right)^2+(1-p_{m,a})
=
\frac{1}{p_{m,a}}-1.
\]
总变差公式则由支撑分解直接给出：
\[
\|\pi_{m,a}-\sigma_m\|_{\mathrm{TV}}
=
\frac12\sum_{x\in A_m^{\max}}\frac{1-p_{m,a}}{K_m}
+\frac12\sum_{x\notin A_m^{\max}}\pi_{m,a}(x)
=
1-p_{m,a}.
\]

最后，由定理 \ref{thm:xi-fixed-freezing-escort-renyi-spectrum-collapse} 的第四条，
\[
1-p_{m,a}\le \exp\!\bigl(-m\Delta(a)+o(m)\bigr).
\]
当 \(m\) 充分大时有 \(1/p_{m,a}\in[1,2]\)。由于固定的凸函数 \(f\) 在紧区间 \([1,2]\) 上 Lipschitz，故
\[
f\!\left(\frac{1}{p_{m,a}}\right)=O(1-p_{m,a}).
\]
代回一般闭式即得
\[
D_f(\sigma_m\|\pi_{m,a})
=
O(1-p_{m,a})
=
O\!\bigl(\exp(-m\Delta(a)+o(m))\bigr).
\]
证毕。
\end{proof}

\begin{conclusion}[冻结相 \texorpdfstring{$\min$}{min}-entropy 的精确有限尺度修正]\label{con:xi-time-part57ba-freezing-minentropy-finite-size}
沿用结论 \ref{con:xi-time-part57ba-freezing-escort-groundstate-fdiv-closed} 的记号。则 escort 分布的 \(\min\)-entropy 满足精确恒等式
\[
\boxed{\;
H_\infty(\pi_{m,a})
:=
-\log\max_{x\in X_m}\pi_{m,a}(x)
=
\log K_m-\log p_{m,a}.\;}
\]
从而
\[
\boxed{\;
H_\infty(\pi_{m,a})
=
\log K_m
+O\!\bigl(\exp(-m\Delta(a)+o(m))\bigr).\;}
\]
特别地，若 \(g_*=\lim_{m\to\infty}m^{-1}\log K_m\)，则
\[
\boxed{\;
H_\infty(\pi_{m,a})-g_*m
=
\bigl(\log K_m-g_*m\bigr)
+O\!\bigl(\exp(-m\Delta(a)+o(m))\bigr).\;}
\]
\end{conclusion}
\begin{proof}
由上一结论，对 \(x\in A_m^{\max}\) 有
\[
\pi_{m,a}(x)=\frac{p_{m,a}}{K_m},
\]
而对 \(x\notin A_m^{\max}\) 有 \(\pi_{m,a}(x)\le p_{m,a}/K_m\)，因为这些点的权重对应严格小于 \(M_m\) 的纤维。故最大概率恰在 \(A_m^{\max}\) 上达到，并满足
\[
\max_{x\in X_m}\pi_{m,a}(x)=\frac{p_{m,a}}{K_m}.
\]
取负对数即得第一式。

再由
\[
1-p_{m,a}\le \exp\!\bigl(-m\Delta(a)+o(m)\bigr)
\]
与标准估计
\[
-\log p_{m,a}=O(1-p_{m,a}),
\]
便得到
\[
H_\infty(\pi_{m,a})
=
\log K_m+O\!\bigl(\exp(-m\Delta(a)+o(m))\bigr).
\]
最后一式只是把主项 \(g_*m\) 从 \(\log K_m\) 中分离出来。证毕。
\end{proof}

\begin{conclusion}[原子 residue 修正的 delta--flat 极值结构]\label{con:xi-time-part57ba-atomic-residue-delta-flat-extremizer}
沿用定理 \ref{thm:xi-atomic-witt-cycle-primitive-exact-peeling} 与定理 \ref{thm:xi-atomic-witt-energy-residue-l2-law} 的记号。固定偶整数 \(n\) 与模数 \(m\ge 2\)，定义原子 residue 修正
\[
A_{m,r}^{\mathrm{cycle}}(n)
:=
A_{m,r}^{\mathrm{full}}(n)-A_{m,r}^{\mathrm{core}}(n),
\]
以及其 twisted Fourier 向量
\[
a_n^{\mathrm{cycle}}(\omega_m^j)
:=
a_n(\omega_m^j)-a_n^{\mathrm{core}}(\omega_m^j).
\]
则有显式公式
\[
\boxed{\;
A_{m,r}^{\mathrm{cycle}}(n)
=
2\,\ind_{\{r\equiv n/2\;(\mathrm{mod}\,m)\}},\;}
\]
\[
\boxed{\;
a_n^{\mathrm{cycle}}(\omega_m^j)=2\,\omega_m^{j n/2}\qquad(j\in\ZZ/m\ZZ).\;}
\]
因此：
\begin{enumerate}
  \item residue 侧支撑恰为单点，即
  \[
  \boxed{\;\#\supp A_{m,\bullet}^{\mathrm{cycle}}(n)=1.\;}
  \]
  \item twisted Fourier 侧全部模长相同，即
  \[
  \boxed{\;\abs{a_n^{\mathrm{cycle}}(\omega_m^j)}=2\qquad(\forall j\in\ZZ/m\ZZ).\;}
  \]
  \item Parseval 能量满足
  \[
  \boxed{\;
  \sum_{r\in\ZZ/m\ZZ}\abs{A_{m,r}^{\mathrm{cycle}}(n)}^2=4,\qquad
  \frac1m\sum_{j=0}^{m-1}\abs{a_n^{\mathrm{cycle}}(\omega_m^j)}^2=4.\;}
  \]
  \item 对任意 \(1\le p\le \infty\)，都有
  \[
  \boxed{\;\|A_{m,\bullet}^{\mathrm{cycle}}(n)\|_{\ell^p}=2,\;}
  \]
  以及
  \[
  \boxed{\;
  \|a_{m,\bullet}^{\mathrm{cycle}}(n)\|_{\ell^p}
  =
  2m^{1/p}\quad(1\le p<\infty),\qquad
  \|a_{m,\bullet}^{\mathrm{cycle}}(n)\|_{\ell^\infty}=2,\;}
  \]
  其中
  \(
  a_{m,\bullet}^{\mathrm{cycle}}(n):=
  (a_n^{\mathrm{cycle}}(\omega_m^j))_{j\in\ZZ/m\ZZ}.
  \)
\end{enumerate}
换言之，原子 residue 修正在 residue 侧呈现单点支撑，而在 twisted Fourier 侧呈现等幅相位向量；二者由有限循环群上的精确 Fourier 对偶所连接。
\end{conclusion}
\begin{proof}
上述两条显式公式正是定理 \ref{thm:xi-atomic-witt-energy-residue-l2-law} 与定理 \ref{thm:xi-atomic-witt-cycle-primitive-exact-peeling} 的直接重述。由 \(A_{m,r}^{\mathrm{cycle}}(n)\) 只在唯一的剩余类 \(r\equiv n/2\pmod m\) 上取值 \(2\)，立得单点支撑与
\[
\|A_{m,\bullet}^{\mathrm{cycle}}(n)\|_{\ell^p}=2
\qquad(1\le p\le \infty).
\]
另一方面，对每个 \(j\) 都有 \(\abs{\omega_m^{j n/2}}=1\)，故
\[
\abs{a_n^{\mathrm{cycle}}(\omega_m^j)}=2
\qquad(\forall j\in\ZZ/m\ZZ),
\]
从而
\[
\frac1m\sum_{j=0}^{m-1}\abs{a_n^{\mathrm{cycle}}(\omega_m^j)}^2
=
\frac1m\sum_{j=0}^{m-1}4
=
4
\]
以及
\[
\|a_{m,\bullet}^{\mathrm{cycle}}(n)\|_{\ell^p}
=
\left(\sum_{j=0}^{m-1}2^p\right)^{1/p}
=
2m^{1/p}
\qquad(1\le p<\infty),
\]
\[
\|a_{m,\bullet}^{\mathrm{cycle}}(n)\|_{\ell^\infty}=2.
\]
最后，residue 侧平方和只有唯一非零坐标贡献，故
\[
\sum_{r\in\ZZ/m\ZZ}\abs{A_{m,r}^{\mathrm{cycle}}(n)}^2=4.
\]
证毕。
\end{proof}

\begin{conclusion}[full/core residue 差的单维秩一扰动与统一投影界]\label{con:xi-time-part57ba-core-full-residue-rankone-perturbation}
沿用结论 \ref{con:xi-time-part57ba-atomic-residue-delta-flat-extremizer} 的记号，并定义差向量
\[
\Delta A_{m,r}(n)
:=
A_{m,r}^{\mathrm{full}}(n)-A_{m,r}^{\mathrm{core}}(n)
=
A_{m,r}^{\mathrm{cycle}}(n).
\]
若记 \(r_0\equiv n/2\pmod m\)，则 residue 侧有
\[
\boxed{\;
\Delta A_{m,\bullet}(n)=2e_{r_0},\;}
\]
其中 \(e_{r_0}\) 为 \(\CC^{\ZZ/m\ZZ}\) 的标准基向量；而 twisted Fourier 侧则有
\[
\boxed{\;
\bigl(a_n(\omega_m^j)-a_n^{\mathrm{core}}(\omega_m^j)\bigr)_{j\in\ZZ/m\ZZ}
=
2(\omega_m^{j n/2})_{j\in\ZZ/m\ZZ}.\;}
\]
因此 full/core residue 差始终局域在一维坐标子空间中。特别地，对任意满足
\[
Q=Q^*=Q^2,
\qquad
Q\mathbf 1=0
\]
的 Hermitian 投影 \(Q\)，都有统一界
\[
\boxed{\;
\|Q\Delta A_{m,\bullet}(n)\|_2\le 2.\;}
\]
故任意由 residue 向量经过固定正交投影所得到的去均值统计量，其 full/core 差都至多表现为一个单维扰动，并且该扰动的 \(\ell^2\) 幅值与模数 \(m\) 无关。
\end{conclusion}
\begin{proof}
由上一结论，
\[
\Delta A_{m,r}(n)=2\,\ind_{\{r=r_0\}},
\]
故 \(\Delta A_{m,\bullet}(n)=2e_{r_0}\)，这说明 residue 差确实落在一维坐标子空间中。其 twisted Fourier 图像同样由上一结论的显式公式直接给出：
\[
a_n(\omega_m^j)-a_n^{\mathrm{core}}(\omega_m^j)=2\omega_m^{j n/2}.
\]

再由 \(Q=Q^*=Q^2\) 可知 \(Q\) 是 \(\ell^2\) 上的正交投影，从而是压缩算子：
\[
\|Qv\|_2\le \|v\|_2
\qquad(\forall v\in\CC^{\ZZ/m\ZZ}).
\]
取 \(v=\Delta A_{m,\bullet}(n)\)，并利用
\[
\|\Delta A_{m,\bullet}(n)\|_2=\|2e_{r_0}\|_2=2,
\]
即可得到
\[
\|Q\Delta A_{m,\bullet}(n)\|_2\le 2.
\]
证毕。
\end{proof}

\endinput
